\sectioncentered*{Заключение}
\addcontentsline{toc}{section}{Заключение}

В рамках курсового проекта проанализирована предметная область интеллектуальных систем управления знаниями и определены требования к пользовательскому взаимодействию с графовым решателем задач. Рассмотрены сценарии пополнения базы документов, навигации по иерархии «сообщество --- домен --- концепция --- карточка», получения объяснимого ответа и перехода к первоисточнику. Установлено, что пользовательская подсистема должна представлять не только конечный результат, но также состояния фоновой обработки, ошибки и доказательный путь, не раскрывая внутреннюю сложность агентного конвейера.

Спроектирована клиентская архитектура, отделённая от агентного ядра и хранилищ типизированным REST API. Определены состояния авторизации, библиотеки, загрузки, карточки, графа и диалога, а также форматы документов, заданий, графовых элементов, источников и путей рассуждения. Разработаны правила адаптивной компоновки и детерминированный алгоритм размещения сообществ, доменов и концепций в SVG-сцене.

Реализовано одностраничное веб-приложение на React и TypeScript. Созданы экран входа и регистрации, библиотека документов, зона одиночной и пакетной загрузки, отображение этапов обработки, повтор ошибочного задания, скачивание оригинала и удаление документа. Реализовано модальное представление карточки знания с аннотацией, ключевыми положениями, доменом, сообществом и концепциями.

Реализована интерактивная визуализация графа знаний. Сообщества представлены отдельными областями, домены --- опорными узлами, а концепции --- связанными интерактивными элементами. Поддерживаются панорамирование, изменение масштаба, сброс viewport, выделение непосредственных связей и отложенная загрузка карточек выбранной концепции. Данные карточек кэшируются в состоянии экрана, поэтому повторное обращение не создаёт лишний запрос.

Реализован диалоговый экран работы с базой знаний. Вопрос проверяется на клиенте и передаётся серверному агенту-ответчику, после чего результат разделяется на текст ответа, подтверждающие источники и пути по концепциям. Для каждого источника доступно авторизованное скачивание оригинала, а для каждого пути отображаются последовательность концепций и оценка релевантности.

Разработан единый модуль взаимодействия с API. Он выполняет JSON- и multipart-запросы, хранит access-токен в памяти, восстанавливает сессию через refresh-маршрут, объединяет одновременные попытки обновления токена и повторяет исходный запрос после ответа 401. Реализованы преобразование серверных ошибок в понятные сообщения, безопасное скачивание бинарных файлов и частичное обновление нескольких активных заданий при помощи периодического опроса.

Программная реализация протестирована статической проверкой TypeScript, производственной сборкой Vite и ручным прохождением основных и ошибочных пользовательских сценариев. Проверены вход и регистрация, загрузка одного и нескольких документов, ограничения формата и количества файлов, обновление статусов, повтор обработки, просмотр карточки, удаление, навигация по графу, отправка вопроса, восстановление сессии и адаптивная компоновка. Контрольная команда сборки завершена без ошибок.

Оценено качество клиентского приложения. Реализованные компоненты покрывают обязательный пользовательский цикл и сохраняют единое поведение сетевых ошибок. Детерминированная раскладка обеспечивает стабильность представления, а опрос только активных заданий, загрузка карточек по требованию и локальный кэш уменьшают число обращений к API. Семантическая разметка, названия кнопок, видимый фокус и учёт reduced motion формируют базовый уровень доступности.

Личный вклад автора заключается в анализе пользовательских сценариев, проектировании клиентской архитектуры и визуального представления базы знаний, реализации React-компонентов, типизированных API-контрактов, загрузки и библиотеки документов, карточки знания, интерактивного графа, диалогового экрана, адаптивных стилей и обработки ошибок, а также в подготовке и прохождении сценариев тестирования и демонстрации. Реализация агентов, серверного API, баз данных, очереди и объектного хранилища использована через согласованные контракты и не отнесена к индивидуальному вкладу автора.

Ожидаемый эффект применения разработанной подсистемы состоит в снижении порога работы со сложной графовой базой знаний: пользователь взаимодействует с документами, концепциями и доказательными путями через единое визуальное пространство. Дальнейшее развитие предлагается направить на фильтрацию и поиск по графу, подсветку пути ответа непосредственно на карте, виртуализацию больших графов, потоковое обновление статусов, автоматические компонентные и сквозные тесты, а также полное управление клавиатурным фокусом и альтернативное табличное представление графа.
